\documentclass[10pt,a4paper]{article}
\usepackage[utf8]{inputenc}
\usepackage{amsmath}
\usepackage{amsfonts}
\usepackage{amssymb}
\usepackage{graphicx}
\usepackage{hyperref}
\usepackage{natbib}
\usepackage{url}
\usepackage{xcolor}
\usepackage{listings}
\usepackage{setspace}
\usepackage{geometry}
\geometry{a4paper, margin=1in}

\hypersetup{
    colorlinks=true,
    linkcolor=blue,
    filecolor=blue,      
    urlcolor=blue,
    citecolor=blue,
}

\title{BT2C: A Pure Cryptocurrency Implementation with Reputation-Based Proof of Stake}
\author{BT2C Core Development Team}
\date{March 19, 2025}

\begin{document}

\maketitle

\begin{abstract}
This paper introduces BT2C (bit2coin), a cryptocurrency designed to function as a pure medium of exchange and store of value without the overhead of smart contracts or decentralized applications. We present a novel reputation-based Proof of Stake (rPoS) consensus mechanism that addresses the energy consumption concerns of Proof of Work while maintaining security properties through cryptographic verification and economic incentives. BT2C implements a deterministic issuance schedule with a fixed maximum supply, combined with a flexible staking model that optimizes for network security and validator participation. This paper outlines the technical architecture, cryptographic foundations, consensus rules, and economic model of the BT2C network.
\end{abstract}

\section{Introduction}
Cryptocurrencies have evolved significantly since the introduction of Bitcoin \cite{nakamoto2008bitcoin} in 2009. While many projects have expanded into complex platforms supporting smart contracts and decentralized applications, BT2C returns to the original vision of a peer-to-peer electronic cash system with improvements in energy efficiency, transaction finality, and participation accessibility.

The core innovation of BT2C lies in its reputation-based Proof of Stake consensus mechanism, which selects validators based on a combination of stake amount and historical performance metrics. This approach provides three key advantages:

\begin{enumerate}
\item Energy efficiency compared to Proof of Work systems
\item Resistance to centralization through accessible validator requirements
\item Enhanced security through reputation-based incentives
\end{enumerate}

\section{System Architecture}

\subsection{Network Topology}
BT2C implements a peer-to-peer network with specialized validator nodes responsible for block production. The network utilizes a gossip protocol for message propagation with the following components:

\begin{itemize}
\item \textbf{Seed nodes}: Entry points for new validators joining the network
\item \textbf{Validator nodes}: Full nodes with staked BT2C that participate in consensus
\item \textbf{API nodes}: Provide REST and WebSocket interfaces for client applications
\item \textbf{Explorer nodes}: Index and serve blockchain data for user interfaces
\end{itemize}

Network communication occurs over TCP/IP with mandatory TLS encryption. Each node maintains a connection pool with a configurable number of peers, with priority given to connections with validators having higher reputation scores.

\subsection{Data Structures}

\subsubsection{Blocks}
Each block in the BT2C blockchain contains:

\begin{verbatim}
Block {
    height: uint64
    timestamp: uint64
    transactions: Transaction[]
    validator: string (BT2C address)
    reward: float64
    previous_hash: string
    hash: string
}
\end{verbatim}

Block hashes are computed as:
\begin{verbatim}
hash = SHA3-256(height || timestamp || merkle_root || validator || reward || previous_hash)
\end{verbatim}

Where \texttt{merkle\_root} is the root of a Merkle tree constructed from transaction hashes.

\subsubsection{Transactions}
Transactions in BT2C follow a simple structure:

\begin{verbatim}
Transaction {
    sender: string (BT2C address)
    recipient: string (BT2C address)
    amount: float64
    fee: float64
    nonce: uint64
    timestamp: uint64
    signature: string
    hash: string
}
\end{verbatim}

Transaction hashes are computed as:
\begin{verbatim}
hash = SHA3-256(sender || recipient || amount || fee || nonce || timestamp)
\end{verbatim}

The signature is generated using the sender's private key over the transaction hash.

\subsection{Cryptographic Foundations}
BT2C employs multiple cryptographic primitives to ensure security:

\begin{enumerate}
\item \textbf{Key generation}: 2048-bit RSA key pairs
\item \textbf{Address derivation}: 
   \begin{verbatim}
   address = "bt2c_" + base58encode(RIPEMD160(SHA256(public_key)))
   \end{verbatim}
\item \textbf{Transaction signing}: RSA-PSS with SHA-256
\item \textbf{Block and transaction hashing}: SHA3-256
\item \textbf{Seed phrases}: BIP39 with 256-bit entropy (24 words)
\item \textbf{HD wallet derivation}: BIP44 path m/44'/999'/0'/0/n
\end{enumerate}

\section{Consensus Mechanism}

\subsection{Reputation-Based Proof of Stake (rPoS)}
BT2C introduces a reputation-based Proof of Stake consensus mechanism that extends traditional PoS by incorporating historical performance metrics into validator selection. The probability of a validator being selected to produce the next block is determined by:

\begin{equation}
P(v) = \frac{s_v}{S_{total}} \cdot R_v
\end{equation}

Where:
\begin{itemize}
\item $P(v)$ is the probability of validator $v$ being selected
\item $s_v$ is the stake of validator $v$
\item $S_{total}$ is the total stake in the network
\item $R_v$ is the reputation score of validator $v$ (range: 0.1 to 2.0)
\end{itemize}

The reputation score $R_v$ is calculated as:

\begin{equation}
R_v = 0.1 + \min(1.9, (0.4 \cdot A_v + 0.3 \cdot U_v + 0.2 \cdot T_v + 0.1 \cdot H_v))
\end{equation}

Where:
\begin{itemize}
\item $A_v$ is the block validation accuracy (range: 0 to 1)
\item $U_v$ is the uptime score (range: 0 to 1)
\item $T_v$ is the transaction processing efficiency (range: 0 to 1)
\item $H_v$ is the historical participation quality (range: 0 to 1)
\end{itemize}

This formula ensures that even validators with poor reputation maintain a minimum selection probability (10\% of their stake-based probability), while high-performing validators can achieve up to 200\% of their stake-based probability.

\subsection{Block Production}
Block production in BT2C follows a time-based schedule with a target block time of 60 seconds. The block production process:

\begin{enumerate}
\item At each block height, validator selection occurs using the rPoS algorithm
\item The selected validator has 30 seconds to produce and broadcast a valid block
\item If the selected validator fails to produce a block, a new validator is selected
\item This process repeats until a valid block is produced or a timeout occurs
\end{enumerate}

\subsection{Block Validation}
When a validator receives a new block, it performs the following validation steps:

\begin{enumerate}
\item Verify the block structure and hash
\item Verify the selected validator's eligibility
\item Verify each transaction's signature and validity
\item Verify the Merkle root matches the transactions
\item Verify the block reward calculation
\end{enumerate}

A block is considered finalized when it has been built upon by 6 subsequent blocks, providing probabilistic finality similar to Bitcoin but with significantly shorter confirmation times due to the 60-second block interval.

\subsection{Validator States}
Validators in BT2C can exist in four distinct states:

\begin{enumerate}
\item \textbf{Active}: Fully participating in consensus
\item \textbf{Inactive}: Registered but not participating (can reactivate)
\item \textbf{Jailed}: Temporarily suspended for protocol violations
\item \textbf{Tombstoned}: Permanently banned for severe violations
\end{enumerate}

Transitions between states follow specific rules:

\begin{itemize}
\item Active $\rightarrow$ Jailed: Occurs when a validator misses more than 50 blocks in a 1000-block window
\item Jailed $\rightarrow$ Active: Requires a waiting period of 100 blocks and a manual unjail transaction
\item Active $\rightarrow$ Tombstoned: Occurs upon detection of double-signing or other critical security violations
\item Tombstoned $\rightarrow$ Any: Not possible; tombstoned validators must create new validator identities
\end{itemize}

\section{Economic Model}

\subsection{Supply Schedule}
BT2C implements a deterministic issuance schedule with the following parameters:

\begin{itemize}
\item \textbf{Maximum supply}: 21,000,000 BT2C
\item \textbf{Initial block reward}: 21.0 BT2C
\item \textbf{Halving interval}: 4 years (126,144,000 seconds, approximately 2,102,400 blocks)
\item \textbf{Minimum reward}: 0.00000001 BT2C
\end{itemize}

The block reward at any height $h$ is calculated as:

\begin{equation}
reward(h) = \max(21.0 \cdot 0.5^{\lfloor h/2102400 \rfloor}, 0.00000001)
\end{equation}

This creates a disinflationary supply curve similar to Bitcoin but with more predictable issuance due to the consistent block time.

\subsection{Fee Market}
Transaction fees in BT2C are dynamic, based on network congestion:

\begin{equation}
min\_fee = base\_fee \cdot (1 + \alpha \cdot utilization\_ratio)
\end{equation}

Where:
\begin{itemize}
\item $base\_fee$ is the minimum fee under no congestion (0.00001 BT2C)
\item $\alpha$ is a multiplier controlling fee sensitivity (currently set to 5)
\item $utilization\_ratio$ is the ratio of transactions in the mempool to the maximum block capacity
\end{itemize}

This mechanism ensures fees remain low during normal operation but increase during periods of high demand to prioritize transactions efficiently.

\subsection{Staking Economics}
BT2C implements a flexible staking model with the following characteristics:

\begin{itemize}
\item \textbf{Minimum stake}: 1.0 BT2C
\item \textbf{No maximum stake}: Validators can stake any amount above the minimum
\item \textbf{No fixed staking period}: Validators can unstake at any time
\item \textbf{Unstaking queue}: Withdrawals enter a FIFO queue processed over time
\item \textbf{Queue processing rate}: Limited to 1\% of total stake per day
\end{itemize}

The effective annual yield for validators is determined by:

\begin{equation}
APY = \frac{block\_rewards\_per\_year \cdot validator\_selection\_probability}{validator\_stake}
\end{equation}

Where $validator\_selection\_probability$ is influenced by both stake amount and reputation as described in section 3.1.

\subsection{Validator Incentives}
\begin{itemize}
\item \textbf{Developer Node Reward}: 1000 BT2C (one-time reward for first mainnet validator)
\item \textbf{Early Validator Reward}: 1.0 BT2C (for validators joining during distribution period)
\item \textbf{Distribution Period}: 14 days from mainnet launch
\item \textbf{Auto-staking}: All distribution period rewards are automatically staked
\end{itemize}

\section{Security Considerations}

\subsection{Sybil Resistance}
The minimum stake requirement of 1.0 BT2C provides a basic economic barrier against Sybil attacks. Additionally, the reputation system requires consistent participation over time to achieve maximum influence, making it costly to establish multiple high-reputation validators.

\subsection{Nothing-at-Stake Problem}
BT2C addresses the nothing-at-stake problem through a combination of:

\begin{enumerate}
\item \textbf{Slashing conditions}: Validators who sign conflicting blocks lose a portion of their stake
\item \textbf{Reputation penalties}: Double-signing results in immediate reputation reduction to minimum
\item \textbf{Tombstoning}: Severe violations result in permanent exclusion from the validator set
\end{enumerate}

\subsection{Long-Range Attacks}
To mitigate long-range attacks, BT2C implements:

\begin{enumerate}
\item \textbf{Weak subjectivity checkpoints}: Published every 10,000 blocks
\item \textbf{Time-bound validator set changes}: Validator set changes take effect only after a delay
\item \textbf{Social consensus backstop}: Community-recognized canonical chain in case of deep reorganizations
\end{enumerate}

\subsection{Transaction Replay Protection}
Each transaction includes a unique nonce derived from the sender's account state, preventing transaction replay attacks. The nonce is incremented with each transaction, and the network rejects transactions with previously used nonces.

\section{Implementation Details}

\subsection{Core Components}
The BT2C implementation consists of several interconnected components:

\begin{enumerate}
\item \textbf{Consensus engine}: Implements the rPoS algorithm and block validation rules
\item \textbf{Transaction pool}: Manages pending transactions and fee prioritization
\item \textbf{State machine}: Tracks account balances, stakes, and validator metadata
\item \textbf{P2P network}: Handles peer discovery and message propagation
\item \textbf{API server}: Provides external interfaces for clients and services
\end{enumerate}

\subsection{Data Persistence}
BT2C uses a hybrid storage approach:

\begin{enumerate}
\item \textbf{Blockchain data}: Stored in a custom append-only file format optimized for sequential access
\item \textbf{State data}: Maintained in a PostgreSQL database for efficient querying and indexing
\item \textbf{Mempool}: Held in memory with Redis backup for persistence across restarts
\end{enumerate}

\subsection{Performance Optimizations}
Several optimizations enable BT2C to maintain high throughput:

\begin{enumerate}
\item \textbf{Parallel transaction verification}: Multiple CPU cores validate transaction signatures concurrently
\item \textbf{Incremental state updates}: State changes are applied incrementally rather than recomputing
\item \textbf{Bloom filters}: Used to quickly check transaction existence without full lookups
\item \textbf{Connection pooling}: Database connections are pooled for efficient resource utilization
\end{enumerate}

\subsection{Network Parameters}
\begin{itemize}
\item \textbf{Target block time}: 60 seconds
\item \textbf{Maximum block size}: 1MB
\item \textbf{Transaction throughput}: Up to 100 tx/s
\item \textbf{Rate limiting}: 100 requests/minute
\item \textbf{Network ports}: 26656 (P2P), 26660 (metrics)
\end{itemize}

\subsection{Infrastructure Requirements}
\begin{itemize}
\item \textbf{Hardware}: 4 CPU cores, 8GB RAM, 100GB SSD
\item \textbf{Software}: Docker \& Docker Compose
\item \textbf{Database}: PostgreSQL
\item \textbf{Caching}: Redis
\item \textbf{Monitoring}: Prometheus \& Grafana
\end{itemize}

\section{Distribution Period Mechanics}
To bootstrap the network securely, BT2C implemented a 14-day distribution period with special incentives:

\begin{enumerate}
\item \textbf{Early validator reward}: 1.0 BT2C for any validator joining during this period
\item \textbf{Developer node reward}: 1000 BT2C for the first validator (network founder)
\item \textbf{Automatic staking}: All distribution rewards are automatically staked
\end{enumerate}

These mechanics were designed to ensure sufficient initial stake distribution while maintaining security during the critical launch phase.

\section{Conclusion}
BT2C represents a focused approach to cryptocurrency design, combining Bitcoin's economic principles with modern consensus mechanisms. By implementing reputation-based Proof of Stake, BT2C achieves energy efficiency without sacrificing security or decentralization.

The deliberate exclusion of smart contracts and complex programmability allows BT2C to optimize for its core use case: a secure, efficient medium of exchange and store of value. This specialization enables performance optimizations and security hardening that would be challenging in more general-purpose blockchain platforms.

As the network continues to mature beyond its initial distribution period, the reputation system will increasingly reward consistent, high-quality participation, creating a virtuous cycle of improved security and performance.

\begin{thebibliography}{9}
\bibitem{nakamoto2008bitcoin}
S. Nakamoto, ``Bitcoin: A Peer-to-Peer Electronic Cash System,'' 2008.

\bibitem{buterin2014slasher}
V. Buterin, ``Slasher: A Punitive Proof-of-Stake Algorithm,'' 2014.

\bibitem{king2012ppcoin}
S. King and S. Nadal, ``PPCoin: Peer-to-Peer Crypto-Currency with Proof-of-Stake,'' 2012.

\bibitem{kiayias2017ouroboros}
A. Kiayias, A. Russell, B. David, and R. Oliynykov, ``Ouroboros: A Provably Secure Proof-of-Stake Blockchain Protocol,'' 2017.

\bibitem{gilad2017algorand}
Y. Gilad, R. Hemo, S. Micali, G. Vlachos, and N. Zeldovich, ``Algorand: Scaling Byzantine Agreements for Cryptocurrencies,'' 2017.
\end{thebibliography}

\end{document}
